\section{The engine}
\label{app:engine}

A benchmark is a call to \texttt{defineBenchmark}. The required pieces are an id, display metadata, tasks or a \texttt{loadTasks} function, a \texttt{grade} function, a provenance array, and a grading-method string. Task ids must be unique. Each task has kind \texttt{single} or kind \texttt{episode}. The returned score must be finite and inside the unit interval, and the reason must be non-empty. The snippet below is the shape from the engine's own documentation. It is not one of the \suitePackages{} packages.

\begin{verbatim}
import { defineBenchmark, exactMatch } from "benchmark-core";

const benchmark = defineBenchmark({
  id: "benchmark-colors",
  name: "Color recall",
  tagline: "Checks exact color recall.",
  tags: ["science"],
  category: "memory",
  whatItMeasures: "Exact recall of a named color.",
  version: "0.1.0",
  status: "alpha",
  tasks: [{
    id: "sky",
    kind: "single",
    prompt: "Name the sky color.",
    expected: "blue",
  }],
  grade: (task, submission) =>
    exactMatch(task.expected, submission.answer ?? ""),
  provenance: [/* name, kind, licence, howToGet, ... */],
  gradingMethod: "Exact string match.",
});
\end{verbatim}

The runner keeps task order and uses bounded concurrency. Each row stores attempts, latency, score, reason, and usage. Episode rows also store the transcript, the tool log, the final state, and an abort reason when a turn budget, a tool-call budget, or a dollar budget trips. A timeout stops the wait. It cannot force a provider that ignores the abort signal to stop spending.

The 95 percent interval is a seeded bootstrap. The default in \texttt{bootstrapCI} is 2{,}000 resamples. The same scores and the same seed give the same interval. That is the interval in Table~\ref{tab:mock}.

Three adapters ship with the engine. \texttt{mockAdapter} takes a seed and, when the benchmark defined them, good and bad scripts. For a single-answer task the good branch returns the expected string and the bad branch returns a fixed wrong token. The choice is one bit: a FNV-1a hash of the task id, mixed with the seed. For an episode the same bit selects the benchmark's good script or its bad script. The CLI does not grow a separate policy flag. \texttt{scripts/collect.mjs} runs seed 1 and seed 2, which flip that bit, and labels each task from the same rule the adapter uses. \texttt{openrouterAdapter} reads \texttt{OPENROUTER\_API\_KEY}. \texttt{openaiCompatAdapter} reads \texttt{OPENAI\_BASE\_URL} and \texttt{OPENAI\_API\_KEY} unless the caller passes them. The runner does not read keys from the repository.

\texttt{replay} resets an environment and applies a recorded tool log, so an episode grade does not need a model call. The reference environment in the engine is a small order lookup that fails once and then returns a fixed order. It tests the engine and sits outside the suite.

\texttt{createCli}, \texttt{createMcpServer}, and \texttt{createHandler} are the other three surfaces, beside the SDK import. The CLI commands are \texttt{list}, \texttt{show}, \texttt{run}, \texttt{score}, \texttt{card}, \texttt{provenance}, \texttt{index-entry}, \texttt{serve}, and \texttt{mcp}. Every command accepts \texttt{-{}-json}. Exit code 0 is success, 1 is bad usage, and 2 is a runtime failure. The MCP server exposes \texttt{list\_tasks}, \texttt{get\_task}, \texttt{score\_answer}, \texttt{run\_benchmark}, and \texttt{get\_provenance}. Only \texttt{run\_benchmark} calls OpenRouter. The HTTP handler is a fetch-style function, so the same code runs on Node and on a Worker. The routes are \texttt{GET /health}, \texttt{GET /v1/tasks}, \texttt{GET /v1/tasks/:id}, \texttt{POST /v1/score}, \texttt{GET /v1/provenance}, and \texttt{GET /v1/index-entry}. The server binds to \texttt{127.0.0.1} unless told otherwise. \texttt{POST /v1/score} has no authentication. A public deployment has to put its own limits in front, especially if a future grader calls a paid judge.

\texttt{card} writes a 1200 by 630 SVG from the report. It is a string template, with no drawing library.

The built-in graders are exact match, a regular expression, a numeric tolerance, and a small JSON Schema subset: type, enum, required, properties, items, additionalProperties, pattern, minimum, and maximum. \texttt{judgeWithLLM} is the optional judge. Its default prompt asks for a JSON object with a numeric score and a string reason. A package that does not call it never spends a judge call.
